% The figures of this section are declared ahead of its text, in the order of their numbers.
\begin{figure*}
\centering
\includegraphics[width=0.24\textwidth]{sequence_0.png}\hfill
\includegraphics[width=0.24\textwidth]{sequence_1.png}\hfill
\includegraphics[width=0.24\textwidth]{sequence_2.png}\hfill
\includegraphics[width=0.24\textwidth]{sequence_3.png}
\caption{Four frames of a collocation trajectory for the mixed cable configuration (B) at $v_0=4\,\mathrm{m/s}$ from $h_0=0.56\Dres$. The quadrotors reposition their cables behind the payload and enter $\Xrf$. The frames show one example of a safe trajectory; the kernel was not computed exhaustively.}\label{fig:sequence}
\vspace{-5pt}
\end{figure*}

\begin{figure}[t]
\centering
\includegraphics[width=0.95\columnwidth]{exact_comparison.png}
\caption{Exact taut-cable model, one initial state of the certified set $\mathcal S$ and one adversarial nominal input, in the plane of wall distance and approach speed; the wall is the right edge. The HOCBF filters lose feasibility (crosses) while every cable leans toward the wall, and reach the wall (triangles). The authority filter brings every cable into the braking half-space (circle) and stops at the wall (square); the dashed curve shows $\Dres$ of its states.}\label{fig:exact}
\end{figure}

\begin{figure}[t]
\centering
\includegraphics[width=0.82\columnwidth]{sampled_margin_map.png}
\caption{Sampled-filter runs in the $(v_0,H(x_0))$ plane (filled: wall contact; open: no contact). The step gives the largest observed contact margin in each speed bin. The observed threshold is not a certified sampled-data bound.}\label{fig:sampled}
\end{figure}

\begin{figure}[t]
\centering
\includegraphics[width=0.90\columnwidth]{sandwich_a.png}
\caption{Viability band for configuration A, cables initially leaning toward the wall: reserved (dashed) and optimistic (dotted) stopping distances. Dots are the smallest \emph{found and re-simulated} collocation trajectories; they do not prove where the kernel boundary lies. Colored curves are two closed-loop trajectories of the sampled authority filter.}\label{fig:sandwich-a}
\vspace{-5pt}
\end{figure}

\section{Numerical study}\label{sec:sim}
Figure~\ref{fig:profile} illustrates the predicted swing profiles. The team has four $1.5\,\mathrm{kg}$ quadrotors carrying a $1\,\mathrm{kg}$ point payload on $1\,\mathrm m$ cables, with $T_{\min}=1\,\mathrm N$, $f_{\max}=44\,\mathrm N$, $\bar T_i=3.2\,\mathrm N$, $\rho_i=34\,\mathrm N$, $a_{\max}=13\,\mathrm{m/s^2}$, $\theta_q=20^\circ$, $\bar z=0.239$, $\bar w=0.171$, $\bar\omega=1\,\mathrm{s^{-1}}$, and $\nu=\nu_w=0.5\,\mathrm{s^{-2}}$. The thrust-reserve bound holds: $\sqrt{(\bar T_i+m_ia_{\max}+m_il_i\bar\omega^2)^2+\rho_i^2}=41.73\,\mathrm N<f_{\max}$. The code and trial videos are available online.\footnote{\url{https://github.com/Hadi-Hajieghrary/Authority_Barriers}}

\textcolor{red}{\textbf{SUGGESTION:} Add a compact experiment-configuration table specifying the $\gamma$ functions and gains, initial-state sampling procedure and seeds, integration settings, solver tolerances, and code version.  Report computation times for the complete filter in Equation~(17), including its altitude constraints.  Specify the numerical backup-CBF baseline's integration horizon and sensitivity-calculation method.  Use comparable accuracy and common hardware for the timing comparison.  State the software environment and number of evaluations.  Report typical and maximum observed computation times and compare them with the stated controller update periods.}

\textcolor{red}{\textbf{WHY:} These settings are needed to reproduce the reported feasibility and wall-contact results.  The comparison of $129$~s against $6$~ms requires enough implementation detail to distinguish an analytical advantage from differences in numerical accuracy or baseline implementation.  The existing wall-only timing measurements do not establish the computational cost of the complete certified filter.  Reporting timing variability also helps reviewers assess whether the implementation can support the proposed update periods.}

\emph{Exact model.} Theorem~\ref{thm:sandwich}(i) was tested in its own setting: the authority filter \eqref{eq:filter} on the taut-cable model under a zero-order hold, with $T^{\mathrm h}_i=2.7\,\mathrm N$, $\nu'=0.4\,\mathrm{s^{-2}}$, $c_\downarrow=5$ and $c_\uparrow=2\,\mathrm{m/s^2}$, $\epsilon=0.1\,\mathrm s$, $c_a=1.25$, and the band $0\le\chi\le45\,\mathrm m$. The 100 initial states have random swing states, speeds of $1$--$4\,\mathrm{m/s}$, and $H(x_0)\in[0,0.05\,\Dres(x_0)]$, between $1\,\mathrm{mm}$ and $0.38\,\mathrm m$; the nominal input is adversarial: the full thrust of every quadrotor toward the wall, tilted $40^\circ$ upward; a trial lasts $6\,\mathrm s$. Table~\ref{tab:exact} gives the result. With a period of $1\,\mathrm{ms}$ the barrier stayed above $-0.05\,\mathrm{mm}$, one trial reached the wall, by $2\,\mu\mathrm m$, and the altitude stayed in its band in every trial; the payload climbed by up to $29.9\,\mathrm m$. 

\textcolor{red}{SUGGESTION: Strengthen the numerical verification of the complete filter in Equation~(17), including its altitude and swing constraints.  Repeat identical initial states and nominal inputs with command holds of $1$, $0.5$, and $0.25$~ms.  At each fixed hold, independently refine the integration settings and solver tolerances.  Report minimum wall clearance, minimum barrier values, and any input or operational-constraint violations.  Evaluate these quantities between controller updates.  Determine how the reported $2,\mu\mathrm{m}$ wall penetration changes under refinement.  Clearly distinguish the continuous-time guarantee from its sampled implementation.}

\textcolor{red}{WHY: Theorem~9(i) provides a continuous-time guarantee, while Table~I evaluates commands applied through a zero-order hold.  The small wall penetration could reflect sampling, integration error, or optimization tolerances.  Independent refinement is needed to distinguish these effects.  The full-order wall-only trials omit constraints from the certified filter and introduce slack-cable dynamics.  Those trials demonstrate limitations but cannot verify the complete certificate under its stated assumptions.}

Figure~\ref{fig:exact} compares the authority filter with HOCBF filters with linear gains $(k_1,k_2)\in\{(0.5,10),(1,20),(2,40)\}\,\mathrm{s^{-1}}$ from one state with $v=3\,\mathrm{m/s}$, every cable at $z_i=-0.8\bar z$, and $h=1.02\,\Dres=7.76\,\mathrm m$. Each HOCBF program is feasible at the start and infeasible after $0.15$, $0.52$, and $0.78\,\mathrm s$, with every cable leaning toward the wall; its relaxation with a penalized slack reaches the wall at $7.3$--$9.2\,\mathrm{m/s}$. The authority filter has every cable in the braking half-space after $0.81\,\mathrm s$ and stops the payload $10\,\mathrm{mm}$ from the wall, with $H$ between $0.15$ and $0.01\,\mathrm m$. Without the bound on the swing accelerations, the filter follows the nominal input with swing accelerations of about $20\,\mathrm{s^{-2}}$, forty times $\nu$; the stopping distance then changes within one period by more than its gradient predicts, and the contacts multiply; a margin of the kind derived for sampled-data CBFs~\cite{breeden2022sampled} is not available for $\Dres$. With the hover floor the stopping distance from $3\,\mathrm{m/s}$ with every cable at $-\bar z$ is $9.18\,\mathrm m$, against $6.56\,\mathrm m$ with the floor $T_{\min}$ and the time-optimal swing: carrying the weight costs $2.6\,\mathrm m$.

\begin{table}[t]
\caption{The authority filter on the exact taut-cable model, 100 trials per row.}\label{tab:exact}
\centering\footnotesize\setlength{\tabcolsep}{4pt}
\begin{tabular}{@{}cccccc@{}}
\toprule
Period & $|\ddot\zeta|\le c_a\nu_\zeta$ & Contacts & Deepest & $\min H$ & $\chi$ in band\\
\midrule
$1\,\mathrm{ms}$ & yes & 1 & $0.002\,\mathrm{mm}$ & $-0.05\,\mathrm{mm}$ & 100\\
$5\,\mathrm{ms}$ & yes & 4 & $0.4\,\mathrm{mm}$ & $-7\,\mathrm{mm}$ & 100\\
$1\,\mathrm{ms}$ & no & 6 & $0.1\,\mathrm{mm}$ & $-13\,\mathrm{mm}$ & 97\\
$5\,\mathrm{ms}$ & no & 33 & $1.7\,\mathrm{mm}$ & $-104\,\mathrm{mm}$ & 100\\
\bottomrule
\end{tabular}
\end{table}

\emph{Full-order model.} The remaining trials use rigid quadrotors, spring--damper cables that can go slack, and a filter at $200\,\mathrm{Hz}$. They were run with the wall barrier alone, that is, with the floor $T_{\min}$ in $\phi_i$, the time-optimal swing, and neither \eqref{eq:filter-alt} nor the bound on the swing accelerations; the altitude is free in them, and the payload climbs by $54\,\mathrm m$ in one trial before it stops. All 690 HOCBF trials lost QP feasibility before a cable entered the braking half-space: the three gain pairs of Fig.~\ref{fig:exact}, each from 100 states with every cable leaning toward the wall at $2$--$4\,\mathrm{m/s}$, with thrust vectors applied directly and with a $1\,\mathrm{kHz}$ attitude loop, and from 30 states with three quadrotors. This does not establish initial feasibility for arbitrary gains. Of 100 trials started with $H(x_0)\in[0,0.05\,\Dres]$ under the adversarial nominal input, the sampled authority filter reached the wall in 59 despite a feasible QP at its sampling instants (Fig.~\ref{fig:sampled}); in 56 trials a cable was slack at some sample. The largest barrier loss was $0.85\,\mathrm m$ within one $5\,\mathrm{ms}$ hold; no contact was observed above an initial margin of $0.78\,\mathrm m$ in a second sweep of 100 states. This margin is empirical and holds for the tested distribution only. The attitude loop incurred $17$--$43\,\mathrm N$ tracking errors during command jumps.

\textcolor{red}{\textbf{SUGGESTION:} Repeat the collocation comparison using the same taut-cable dynamics, vehicle count, input limits, operational constraints, and altitude band as the certified filter.  Match all initial-state components except the wall clearance being minimized.  Specify the tension floors and swing maneuver used to compute $D$ in each comparison.  Verify the hover, thrust, and swing reserve assumptions for each vehicle configuration.  Re-simulate candidate trajectories and check every constraint.  Report the smallest found safe clearance relative to $D$.  Include a representative tighter altitude band to assess practical usefulness.}

\textcolor{red}{\textbf{WHY:} The reported clearances of $0.25D$ to $0.55D$ come from searches with different model and constraint settings from the complete certificate.  Relaxed constraints can permit shorter stopping distances.  A matched comparison is needed to determine how much of the gap comes from the certificate's conservative construction.  Testing a tighter altitude band would also clarify its operating usefulness, given the reported climb of up to $29.9$~m.}

\emph{Conservatism.} Direct-collocation searches with three quadrotors covered three initial cable configurations at rest (A: $z_i=-0.15$; B: $z=(-0.15,0,0.15)$; C: $z_i=\bar z$) at eight speeds from $0.5$ to $4\,\mathrm{m/s}$, with each candidate subsequently simulated; like the full-order trials, they constrain the wall distance only. When cables already lean away from the wall, the observed feasible distance is about $0.97$--$0.98$ times $\Dres$; when a swing is required, feasible trajectories can begin at only $0.25$--$0.55\,\Dres$ (Fig.~\ref{fig:sandwich-a} shows configuration A). Figure~\ref{fig:sequence} shows one such trajectory. The large gap is driven in part by the reserved swing acceleration: $\nu=0.5\,\mathrm{s^{-2}}$ needs $1.96\,\mathrm{s}$ to cross the allowed angle range, while the physical thrust can reach the swing-rate limit much sooner, and Assumption~\ref{ass:op} admits no $\nu$ above $0.54\,\mathrm{s^{-2}}$ at this rate limit. A collocation success gives an example trajectory, but a failure does not locate the viability boundary, so the evidence supports the qualitative gap between the two models and validates neither set inclusion.

\emph{Cost.} A backup filter that integrates $\plan$ numerically \cite{chen2021backup} reproduced $\Dres$ within a relative error of $1.1\cdot10^{-6}$ and the commands of the filter within $0.4\,\%$ of $f_{\max}$ at 60 states, at $129\,\mathrm s$ per evaluation against $6\,\mathrm{ms}$. In Python with the solver Clarabel, evaluating $\Dres$ takes $0.5$--$0.6\,\mathrm{ms}$ and the filter with the wall barrier $2$--$4.5\,\mathrm{ms}$ for two to eight quadrotors, close to the $5\,\mathrm{ms}$ sampling period; solver failures, model mismatch, and inter-sample excursions preclude a real-time guarantee.
